\section{Results}
\label{sec:results}

\subsection{Support at most two is exact under the unit objective}
\label{sec:support-theorem}

\begin{theorem}[Support reduction]
For every qubit count, target six-tuple, matching, relative target assignment, and central-slot choice in the grammar of Sec.~\ref{sec:setting}, the unit-objective optimum is attained by a configuration in which every frame string has support at most two. Hence
\[
\cstar(t)=\ctwo(t).
\]
\end{theorem}

The proof is local but yields an all-size statement. For a frame $R$ with support $w\geq3$ and anticommuting partner $R'$, label each support coordinate by
\[
(\alpha_q,\beta_q)=
(\langle R_q,R'_q\rangle,\langle S_q,R_q\rangle)\in\Ftwo^2.
\]
Anticommutation makes the total $\alpha$ parity odd. Any such multiset of size at least three contains a nonempty proper singleton or pair whose $\alpha$ and $\beta$ sums both vanish. Zeroing those letters therefore preserves the relevant anticommutation and tag syndrome without tag repair. A complete 18,432-case local enumeration verifies that the increase of the restoration term in Eq.~\eqref{eq:cost} never exceeds the refunded frame cost; there are zero violations. Repeating the exchange decreases total frame support without increasing cost. Lexicographic descent on cost and total frame support terminates at support at most two. Appendix~\ref{app:support-proof} gives the argument and its exact computational obligation.

Support two is a genuine boundary of this proof. At $w=2$, four odd-parity class patterns lack a nonempty proper zero-sum subset. Thus the theorem does not follow by simply repeating a support-one argument.

\subsection{Finite evidence refutes the initial formula}

Table~\ref{tab:initial} separates the benchmark components. The initial formula matches the complete structured two-qubit slice and all application-derived rows. In the seeded panel it misses exactly one three-qubit input. Across the 9,546 distinct benchmark inputs, the outcome is therefore 9,545 matches and one mismatch.

\begin{table}[t]
\centering
\caption{Exact comparison of the initial explanatory formula. The seeded panel includes the retained refuting input.}
\label{tab:initial}
\begin{tabular}{lrrr}
\toprule
Panel & Distinct inputs & Matches & Mismatches\\
\midrule
Structured two-qubit slice & 9,261 & 9,261 & 0\\
Fresh seeded two-/three-qubit panel & 240 & 239 & 1\\
Application-derived subjects & 45 & 45 & 0\\
\midrule
Total & 9,546 & 9,545 & 1\\
\bottomrule
\end{tabular}
\end{table}

The mismatch has $n=3$ and target pairs
\[
((XYZ,YYX),(YYX,XXZ),(ZZI,IYI)).
\]
Exact optimization gives $\cstar=10$, whereas all three terms in Eq.~\eqref{eq:simple} equal 11. Thus
\begin{equation}
\cstar=10<11=\csimple.
\label{eq:counterexample}
\end{equation}
Its minimizing support-two frame uses a borrow home outside its block's own target support. The row refutes only the completeness of the three-family explanation. It does not refute feasibility of the formula, because $10\leq11$, or the support theorem, because $\cstar=\ctwo=10$.

Before this mismatch was known, 102 predictions had been staged and then exactly evaluated: 90 public-library matchings and 12 engineered boundary rows, with four examples allocated to each of the three initial configurations. All 102 matched. These prospective successes provide calibration evidence, but Eq.~\eqref{eq:counterexample} shows why they cannot establish global equality.

\subsection{The exact evaluator passes its implementation checks}

The support-two evaluator returns the exact cost on all 9,547 comparison events, representing 9,546 distinct inputs. It also returns 10 on the refuting row. The unrestricted optimizer is absent from its forecast path. For higher-qubit application rows, exactness is sometimes certified by the containment pinch
\[
\cstar\leq\ctwo\leq C_{\mathrm{wa}}=\cstar,
\]
rather than by a direct support-two sweep. This distinction matters: zero finite error validates consistency of the code and witnesses, whereas the support-reduction proof is the source of the all-size claim.

\subsection{Explanation repair, hostile refutation, and all-size closure}

Allowing borrow homes outside a block's own target support repairs Eq.~\eqref{eq:counterexample}. That repair is still incomplete. A frozen hostile search evaluates 740 exact cases and finds 64 in which the enlarged-borrow explanation remains one unit above $\ctwo$. All 64 optima stay within support two, so the exact forecast is unaffected.

Adding the hybrid family covers 10,481 verified finite inputs, including the 64 hostile witnesses, with no observed need for a fifth configuration. This finite zero-residual result is useful but is not the all-size proof. The final proof isolates the remaining commuting support-two blocks and admits independent permutations of each block's two targets, a degree of freedom already present in the feasible configuration space. The complete domain comprises 378 irreducible geometries and $16{,}777{,}216$ target-letter states per geometry. Every state has a non-increasing alternative with strictly fewer commuting support-two blocks; no residue remains. Induction on their number proves Eq.~\eqref{eq:classification} for all qubit counts under the unit objective.

\begin{table}[t]
\centering
\caption{Certificate chronology. A finite check is not promoted into the proof column.}
\label{tab:chronology}
\begin{tabularx}{\textwidth}{>{\raggedright\arraybackslash}p{.23\textwidth}>{\raggedright\arraybackslash}p{.27\textwidth}>{\raggedright\arraybackslash}X}
\toprule
Layer & Exact evidence & Conclusion\\
\midrule
Initial explanation & 9,545/9,546 distinct inputs; 102/102 prospectively staged rows & Refuted by one exact row\\
Exact forecast & Support-reduction theorem; 9,547 zero-error comparison events & Exact for all sizes in scope\\
Enlarged borrow & 64 failures among 740 hostile cases & Explanation refuted; cost theorem intact\\
Hybrid explanation & 10,481 finite inputs with no residual & Finite coverage only\\
Final explanation & Complete domains for 378 geometries; no local residue & All-size classification under the unit objective\\
\bottomrule
\end{tabularx}
\end{table}

\subsection{Adverse boundary under a changed objective}
\label{sec:adverse-objective}

A pre-specified reweighted objective changes the local exchange economics. On one exact witness,
\begin{equation}
\cstar=11<\ctwo=13<C_{\mathrm{wa}}=23.
\label{eq:reweighted}
\end{equation}
The corresponding panel contains 53 failures of support-two closure. Equation~\eqref{eq:reweighted} is not a perturbative degradation of the main theorem; it falsifies transfer of the theorem to that objective. Any new weighting requires a fresh exchange inequality or a different structural argument.
